\scanpage{211}
\inserted{Comme $\overline V$ est quasi-compact, on voit que \uline{à condition de} prendre $Y$ assez petit, $H'$ \uline{est fermé dans $P^\tau$}. Même conclusion si $H$ est propre sur $\unsure{Y|S}$.} Il suffit de prouver que \deleted{\unsure{à condition,} \unsure{que}} ces \deleted{parties de $Y$}\inserted{conditions}, on a $H=P'$. En effet, il \deleted{\unsure{contient}} contient les $P_s^0$. Il suffit donc, compte tenu que $P_s^0$ est divisible pour les entiers premiers à la car. de $k(s)$, de prouver que \deleted{les éléments de $P_s$} pour \unsure{tt|tout} entier $n$ premier à la car. de $k(s)$, $ {}_n(P_s)\subset H$. Or $ {}_nP_s=\underline{\mathscr H}^{1}(X_s,{}_n\mu)$, d'autre part on sait \unsure{(réf)|(cf)} que sur l'ouvert \inserted{$S_n$} où $n$ est inversible, $\underline{\mathscr H}^{1}(X_n/S_n,{}_n\mu)$ est un \deleted{\unsure{schéma} \uline{\unsure{vectoriel|en groupes}} de} \uline{schém. étale} sur $S$, dont \unsure{tt|tout} pt de sa fibre en $s$ est spécialisation d'un point de la fibre générique, i.e. d'un pt de $ {}_n(P_{\unsure{g|\eta}})\subset P'_g\subset H$. Comme $H$ est fermé, on a donc $ {}_n(P_s)\subset H$ cqfd

\uline{Corollaire} Supposons que (i) \inserted{les $(P_s)^0$ sont propres ou} dans $P^\tau$ \deleted{\unsure{les}} l'adhérence d'un pt soit quasi-compacte (ii) les $P_s^\tau/P_s^0$, \deleted{\unsure{pour $s\in S$, sont propres sur $k(s)$}} \inserted{sont d'ordre premier à la \unsure{car|c|car.} $k(s)$}. Alors $P^\tau$ est de type fini sur $S$.
% archive sheet number: 210
